\documentclass[10pt,a4paper]{amsart}
\usepackage[margin=19mm]{geometry}
\usepackage{amsmath,amssymb,mathtools}
\usepackage[scr=rsfs,cal=euler]{mathalfa}
\usepackage{xfrac}
\usepackage[unicode,hidelinks,bookmarks=false]{hyperref}
\newcommand{\ie}{i.e.\ }
\newcommand{\cf}{cf.\ }
\newcommand{\resp}{respectively\ }
\newcommand{\Subsection}{\subsection}
\providecommand{\nobreakdash}{\nobreak\hbox{-}}
\setlength{\emergencystretch}{2em}
\multlinegap0pt
\usepackage{bm}
\usepackage{bbm}
\usepackage{dsfont}
\usepackage{xspace}
\usepackage{cancel}
\usepackage{mathtools}
\usepackage{cleveref}
\usepackage{amsthm}
\makeatletter
\@ifundefined{lemma}{\newtheorem{lemma}{Lemma}}{}
\makeatother
\newcommand{\comp}[1]{\bar{#1}}
\newcommand{\wt}{s_p}
\title[A degree bound for planar functions]{A degree bound for planar functions}
\author{Christof Beierle, Tim Beyne}
\date{Source: arXiv:2407.04570v2 (CC BY 4.0)}
\begin{document}
\maketitle

\noindent\emph{Source and licence.} A degree bound for planar functions. Version arXiv:2407.04570v2 (CC BY 4.0), distributed under the Creative Commons Attribution 4.0 International licence, as recorded in the arXiv licence metadata for this exact version and shown on the abstract page https://arxiv.org/abs/2407.04570v2. Full rights notice: licenses/arxiv-2407.04570-CC-BY-4.0.txt.\par\smallskip

\noindent\textit{Editorial scope.} Counted unit: the Lemma whose source label is \texttt{lem:remaining\_p7} in the article (source0.tex lines 554--557, source label \texttt{lem:remaining\_p7}), delivered together with the article's own complete original proof of it (source0.tex lines 558--682, one \texttt{proof} environment of 125 lines). No mathematics is written, repaired or re-derived: statement and proof are the article's own text line for line. Required context retained uncounted: eq:general\_form (lines 442--445); lem:power\_2\_finite (lines 461--463); eq:de (lines 474--477); lem:technical (lines 519--527). Every definition, display and label of those lines is retained unchanged. Omitted from this package and not redistributed: 1--441, 446--460, 464--473, 478--518, 528--553, 683--1028. Nothing inside a retained excerpt is omitted or altered: every retained body is delivered exactly as the article's lines are written, so no omission receipt of any kind accompanies this package. Citations: the retained excerpts carry no citation command at all, so no bibliography entry of the article is transcribed and no reference is redistributed. Reference adaptation: none was needed and none was made; every reference inside the delivered unit resolves inside it, so no unresolved reference or citation is produced. Printed slips kept literal: no printed word, symbol, hypothesis or number of the counted statement or of its proof was altered. Layout and class: the article's own class is \texttt{article} with packages including ct\_modified, enumerate, mathrsfs, thmtools, thm-restate; the wrapper is the lane's pinned \texttt{amsart} editorial wrapper at 10pt on A4, which restates the theorem-like environments the retained text needs and adds the article's own macro definitions that the retained text needs (\texttt{\textbackslash comp}, \texttt{\textbackslash wt}), with extra wrapper packages (bm, bbm, dsfont, xspace, cancel, mathtools, cleveref). The delivered numbering is the unit's own. Assets: no retained span contains a figure environment, a graphics-inclusion command, a PDF-page inclusion, a table or a TikZ picture; no image file is copied into this package, the package declares no asset dependency and no image file is redistributed with it. Licence: version arXiv:2407.04570v2 of this article is distributed under the Creative Commons Attribution 4.0 International licence, as recorded in the arXiv licence metadata for this exact version and shown on the abstract page https://arxiv.org/abs/2407.04570v2. A full rights notice is delivered at licenses/arxiv-2407.04570-CC-BY-4.0.txt. Characters: every retained excerpt is pure ASCII, so the delivered engine, XeLaTeX with its default Latin Modern text font, reports no missing character.
\begin{align}
\label{eq:general_form}
(\overbrace{0,\dots,0}^{r}, t, \overbrace{u_1,\dots,u_s}^{s},1,\overbrace{0,\dots,0}^{r},p-t, \overbrace{\comp{u}_1,\dots,\comp{u}_s}^s,0)_p,
\end{align} 

\begin{lemma} \label{lem:power_2_finite}
Suppose $p>3$, $n = 2m$ with $m \geq 2$ and let $s$ be a non-negative integer. If $d = D(t,u_1,\dots,u_s)$ with $t, u_1,\dots,u_s$ in $\{0,\dots,p-1\}$ and $t\neq 0$ such that $t \notin\{1,p-1\}$ or $u_j \notin \{0,1,p-2,p-1\}$ for some $j$ in $\{1,\dots,s\}$, then there exists an  $e$ in $\{1,\dots,p^n-1\}$ such that $\wt(e \star d) - \wt(e) > m(p-1)$.
\end{lemma}

\begin{align}
\label{eq:de}
    [\overbrace{0,\dots,0}^r,t,\overbrace{u_1,\dots,u_s}^s,2h+1,\overbrace{0,\dots,0}^r,p-t,\overbrace{\comp{u}_1,\dots,\comp{u}_s}^s,2h]_p.
\end{align}

\begin{lemma}
\label{lem:technical}
Let $p > 3$. Let $s$ be a positive integer and $u_0,u_1,\dots,u_s \in \{0,1,p-2,p-1\}$. For $v = [2u_1-u_0,\dots,2u_{s-1}-u_{s-2},2u_s -u_{s-1}]_p \neq p^s-1$, there exists $\delta$ in $\{-1, 0, 1\}$ such that
\begin{equation*}
v + \delta\,p^s =  [\gamma_1,\dots,\gamma_s]_p,
\end{equation*}
with $\gamma_1, \ldots, \gamma_s$ in $\{0, \ldots, p - 1\}$ and $0 < u_s + \delta + 1 \le p - 1$.
Furthermore, if $u_s = 1$, $\delta = 1$ and $p > 5$, then there exists an index $j$ such that $\gamma_j \not\in \{0, 1, p - 2, p - 1\}$.
\end{lemma}

\begin{lemma}
\label{lem:remaining_p7}
Let $p>5$, $n = 2m$ with $m \geq 2$, and $s$ a non-negative integer. If we have $d = D(t,u_1,\dots,u_s)$ with $t$ in $\{1,p-1\}$ and $u_1,\ldots,u_s$ in $\{0,1,p-2,p-1\}$, then there exists an $e$ in $\{1,\dots,p^n-1\}$ such that $\wt(e \star d) - \wt(e) > m(p-1)$.
\end{lemma}

\begin{proof}
Let $r = m - 2 - s$.
We will multiply $d$ by $e = hp^m + h+1 + (p-1)p^m = (p-1+h)p^m + h+1$ for various choices of $h$ with $1 \leq h < p^{m+1}$.
The case $s = 0$ will be handled separately from the case $s > 0$.

\paragraph{Case $s=0$.} 
We first multiply $d$ by $(p-1)p^m$ and obtain (see \eqref{eq:general_form} for $d$)
\begin{equation*}
d (p-1)p^m \equiv [\overbrace{0,\dots,0}^r,t,p-1-t,\overbrace{0,\dots,0}^r,-t,p-1+t]_p \pmod{p^n-1}.
\end{equation*}
Let $e = hp^m + h+1 + (p-1)p^m$. By adding the term \eqref{eq:de} from the proof of \Cref{lem:power_2_finite}, we obtain 
\begin{equation*}
    ed \equiv [\overbrace{0,\dots,0}^r,2t, 2h-t+p,\overbrace{0,\dots,0}^r,p-2t, p-1+2h+t]_p \pmod{p^n-1}.
\end{equation*}
If $t = p-1$, then $2t = p + p-2$, so we can write $e \star d$ (mod $p^n-1$) as
\begin{equation*}
    [\overbrace{0,\dots,0}^r,p-2, 2h+2,\overbrace{0,\dots,0}^r,2, p-2+2h+p-1]_p.
\end{equation*}
In particular, whenever $t \in \{1,p-1\}$, we have
\begin{equation*}
    e \star d \equiv [\overbrace{0,\dots,0}^r,t', 2h-u_0+p,\overbrace{0,\dots,0}^r,p-t', p-1+2h+u_0]_p \pmod {p^n-1},
\end{equation*}
where $t' \in \{2,p-2\}$, and $u_0 = 1$ if $t' = 2$ and $u_0 = p - 2$ if $t' = p - 2$.

If $t' = 2$, let $h = p^{r + 1} - 1 - (p - 1)/2$. As $0 \leq h+1 <p^m$ and $0 \leq p-1+h <p^m$, we have $\wt(e) = \wt(h+1) + \wt(p-1+h)$. Since $\wt(h + 1) = r(p-1) + (p - 1)/2 + 1$ and $\wt(p - 1+ h) =  (p - 1)/2$, it follows that $\wt(e) = (r + 1)(p - 1) + 1$.
From $2h = p^{r + 1} - 1 + \sum_{i = 1}^r (p - 1)p^i$, we obtain
\begin{align*}
e \star d = [&\overbrace{p-1,\dots,p-1}^r,t'+1, p-1-u_0,\overbrace{p-1,\dots,p-1}^r,p-t'+1, p-2+u_0]_p \\
= (&\overbrace{p-1,\dots,p-1}^r,3, p-2,\overbrace{p-1,\dots,p-1}^r,p-1, p-1)_p.
\end{align*}
Hence, $\wt(e \star d) = 2(r+1)(p-1) + p+1 = m(p-1) + (r+1)(p-1) + 2 > m(p - 1) + \wt(e)$.

If $t' = p-2$, let $h = p^{r + 1} - (p - 1)$. We then have $\wt(e) = \wt(h+1) + \wt(p-1+h) = 2 + r(p-1) + 1 = r(p-1) + 3$ and $2h = 2 - p + p^{r+1} + \sum_{i=1}^r(p-1)p^i$. If $r>0$, we obtain
\begin{align*}
    e \star d \equiv [&p-2,\overbrace{p-1,\dots,p-1}^{r-1},t'+1, 2-u_0+p,\\
    &p-2,\overbrace{p-1,\dots,p-1}^{r-1},p-t'+1, p+1+u_0]_p \\
    \equiv [&p-2,\overbrace{p-1,\dots,p-1}^{r-1},p-1, 4,p-2,\overbrace{p-1,\dots,p-1}^{r-1},3, p+p-1]_p \\
    \equiv (&p-1,\overbrace{p-1,\dots,p-1}^{r-1},p-1, 4,p-2,\overbrace{p-1,\dots,p-1}^{r-1},3,p-1)_p \pmod {p^n-1},
\end{align*}
with $\wt(e \star d) = 2(r+1)(p-1)+6 = (r+2)(p-1) + r(p-1)+6 > m(p - 1) + \wt(e)$. If $r = 0$, then we obtain (using $t' = p - 2$ and $u_0 = p - 2$) that $e \star d$ is congruent to
\begin{align*}
    [t', 2-u_0+p,p-t', p+1+u_0]_p &\equiv [p-2, 4,2, p+p-1]_p \\
    &\equiv (p-1, 4,2, p-1)_p \pmod {p^n-1}.
\end{align*}
Hence, $\wt(e \star d) = (r+2)(p-1) + r(p-1) + 6 > m(p - 1) + \wt(e)$.

\paragraph{Case $s>0$.} 
Multiplying $d$ by $(p-1)p^m$ modulo $p^n - 1$ yields
\begin{align*}
[&\overbrace{0,\dots,0}^r,t,u_1-t,\overbrace{u_2-u_1,u_3-u_2,\dots,u_s-u_{s-1}}^{s-1},p-1-u_s,\\
&\overbrace{0,\dots,0}^r,-t,-(u_1-t),\overbrace{-(u_2-u_1),-(u_3-u_2),\dots,-(u_s-u_{s-1})}^{s-1},p-1+u_s]_p.    
\end{align*}
Let $e = hp^m + h+1 + (p-1)p^m$. By addition with term \eqref{eq:de} from the proof of \Cref{lem:power_2_finite}, we obtain that $e \star d$ is congruent to
\begin{align*}
    [&\overbrace{0,\dots,0}^r,\phantom{p-}\;\,2t,2u_1-t,\overbrace{2u_2-u_1,\ldots,2u_s-u_{s-1}}^{s-1}, 2h-u_s+p,\\ &\,0,\dots,0,p-2t, \overline{2u_1-t},\overline{2u_2-u_1},\dots,\overline{2u_s-u_{s-1}},p-1+2h+u_s]_p.
\end{align*}
If $t = p-1$, then we have $2t = p + p-2$, so we can write $e \star d$ (mod $p^n-1$) as
\begin{align*}
    [&\overbrace{0,\dots,0}^r,\hphantom{p-()}p-2,2u_1-(p-2),\overbrace{2u_2-u_1,\dots,2u_s-u_{s-1}}^{s-1}, 2h-u_s+p,\\ &\,0,\dots,0,\,p-(p-2), \overline{2u_1-(p-2)},\overline{2u_2-u_1},\dots,\overline{2u_s-u_{s-1}},p-1+2h+u_s]_p.
\end{align*}
Hence, in any case of $t \in \{1,p-1\}$, we have
\begin{align*}
e \star d \equiv [&\overbrace{0,\dots,0}^r,\hphantom{p-}\;\,t',\overbrace{2u_1-u_0,\ldots,2u_s-u_{s-1}}^{s}, 2h-u_s+p,\\ &0,\dots,0,p-t', \overline{2u_1-u_0},\ldots,\overline{2u_s-u_{s-1}},p-1+2h+u_s]_p \pmod {p^n-1},
\end{align*}
where $t' \in \{2,p-2\}$ and $u_0 = 1$ if $t' = 2$ and $u_0 = p - 2$ if $t' = p - 2$. Define $v = [2u_1-u_0,\dots,2u_s - u_{s-1}]_p$. Note that $v \neq p^s-1$, as $u_0 \neq p-1$.
By \Cref{lem:technical}, there exists a $\delta$ in $\{-1, 0, 1\}$ such that $v = (\gamma_1, \ldots, \gamma_s)_p - \delta p^s$ with $\gamma_1, \ldots, \gamma_s$ in $\{0,\dots,p-1\}$. Hence, $e \star d$ is congruent to
\begin{equation*}
[\overbrace{0,\dots,0}^r,t',\overbrace{\gamma_1,\dots,\gamma_s}^s, 2h-u_s+p-\delta, \overbrace{0,\dots,0}^r,p-t', \overbrace{\comp{\gamma}_1,\dots,\comp{\gamma}_s}^s,p-1+2h+u_s+\delta]_p.
\end{equation*}
We finish the proof using a case distinction between three subcases corresponding to $u_s \notin \{0, 1\}$ and $r = 0$, $u_s \notin \{0, 1\}$ and $r > 0$, and $u_s \in \{0, 1\}$.

\paragraph{Subcase $u_s \notin \{0,1\}$ and $r=0$.} We choose $h = 1$ so that $\wt(e) = 3$. We have
\begin{equation}
\label{eq:fail_p5}
e \star d = [t'+1,\overbrace{\gamma_1,\dots,\gamma_s}^s, p-u_s+2-\delta, p-t', \overbrace{\comp{\gamma}_1,\dots,\comp{\gamma}_s}^s,1+u_s+\delta]_p.
\end{equation}
By assumption, $u_s \in \{p-2,p-1\}$ and $p>5$, so $0 \leq p-u_s+2 - \delta \leq p-1$.  Moreover, by \Cref{lem:technical}, $0 \leq 1+u_s + \delta \leq p-1$. Hence, $\wt(e \star d) = (s+2)(p-1)+6 = m(p-1)+6 > m(p - 1) + \wt(e)$.

\paragraph{Subcase $u_s \notin \{0,1\}$ and $r>0$.}
Let $h = ((p - 1)/2)\,p^r - (p - 1)$. 
Since $\wt(h + 1) = r(p - 1) - (p - 1) / 2 + 1$ and $\wt(p - 1 + h) = (p - 1)/2$, we have $\wt(e) = \wt(h+1) + \wt(p-1+h) = r(p-1) + 1$. Furthermore, $2h = (p - 1) p^r - 2(p - 1) = (p-2)p^r + 2 - p + \sum_{i=1}^{r-1}(p-1)p^i$. Hence, $e \star d$ is
\begin{align*}
[&\overbrace{p-1,\dots,p-1}^{r-1},p-2,\hphantom{p-}\;\,t',\overbrace{\gamma_1,\dots,\gamma_s}^s,2-u_s-\delta, \\ 
&\overbrace{p-1,\dots,p-1}^{r-1},p-2,p-t', \overbrace{\comp{\gamma}_1,\dots,\comp{\gamma}_s}^s,1+u_s+\delta]_p.
\end{align*}
To determine the sum of base-$p$ digits of $e\star d$, we use the same argument as for the case $r=0$ above. If $r-1 > 0$, then $e \star d$ is equal to
\begin{align*}
&(\overbrace{p-1,\dots,p-1}^{r-1},\;\;\;\;\;\;\;\;\,p-2,t',\overbrace{\gamma_1,\dots,\gamma_s}^s, p-u_s+2-\delta,\\ 
&p-2,\overbrace{p-1,\dots,p-1}^{r-2},p-2,p-t', \overbrace{\comp{\gamma}_1,\dots,\comp{\gamma}_s}^s,1+u_s+\delta)_p.
\end{align*}
If $r-1 = 0$, then $e \star d$ equals
\begin{equation*}
(p-2,t',\overbrace{\gamma_1,\dots,\gamma_s}^s,p-u_s+2-\delta, p-3,p-t', \overbrace{\comp{\gamma}_1,\dots,\comp{\gamma}_s}^s,1+u_s+\delta)_p.
\end{equation*}
In both cases, we obtain $\wt(e\star d) = m(p-1) + r(p-1)+2 > m(p - 1) + \wt(e)$. 

\paragraph{Subcase $u_s \in \{0, 1\}$.} Let $z \in \{1,\dots,s+1\}$ and $h = (p^{r + 1 + z} - 1) + (p^{r + 1} - 1) - (p - 1)/2$. Equivalently, $h = p^{r + 1 + z} - 1 + (p - 1)/2 + \sum_{i = 1}^r (p - 1)p^i$.
Let us first analyze the sum of base-$p$ digits of $e$. If $z \neq s+1$, then we have $0 \leq p-1+h <p^m$ and $0 \leq h+1 < p^m$ with $\wt(h+1) = (p-1)/2 + r(p-1) + 1$ and $\wt(p-1+h) = (p - 1)/2$, yielding $\wt(e) = (r+1)(p-1) +1$. If $z = s+1$, then $r+1+z = m$, so $(p-1+h)p^m + h+1$ is congruent to
\begin{align*}
&\frac{p-1}{2} + \sum_{i=1}^r(p-1)p^i + p^m + \left(\frac{p-1}{2}-2\right)p^m + p^{r+1+m} + 1 \\
\equiv &\left(\frac{p-1}{2}+1\right) + \sum_{i=1}^r(p-1)p^i + \left(\frac{p-1}{2}-1\right)p^m + p^{r+1+m} \pmod{p^n - 1}.
\end{align*}
We take $e$ with $1 \leq e \leq p^n-1$ congruent to $(p-1+h)p^m + h+1$, so that
$\wt(e) = (r+1)(p-1) + 1$ as well. 
Multiplying $h$ by $2$, we obtain $2h = p^{r+1} + 2p^{r+1+z} -3 + \sum_{i=1}^r(p-1)p^i$.
We now consider the cases $(u_s,\delta) \neq (1,1)$  and $(u_s,\delta) = (1,1)$ separately.

In case $(u_s,\delta) \neq (1,1)$, choosing $z = s+1$ yields
\begin{align*}
e \star d \equiv [&\overbrace{p-1,\dots,p-1}^r,\hphantom{p-}\;\,t'+1,\overbrace{\gamma_1,\dots,\gamma_s}^s, p-u_s-1-\delta,\\ &p-1,\dots,p-1,p-t'+1, \comp{\gamma}_1,\dots,\comp{\gamma}_s,p-2+u_s+\delta]_p \pmod {p^n-1}. 
\end{align*}
From \Cref{lem:technical}, we have $0 \le u_s + \delta < p - 1$ (note that $v \neq p^s - 1$ as $u_0 \neq p - 1$). Hence, $0 < p - 1 - u_s - \delta \le p - 1$.
Furthermore, by assumption $u_s \in \{0, 1\}$ and $(u_s, \delta) \neq (1, 1)$, so $0 \le p - 2 + u_s + \delta \le p - 1$. Hence,
\begin{align*}
e \star d = (&\overbrace{p-1,\dots,p-1}^r,\phantom{p-}\;\, t'+1,\overbrace{\gamma_1,\dots,\gamma_s}^s, p-u_s-1-\delta,\\ &\,p-1,\dots,p-1,p-t'+1, \comp{\gamma}_1,\dots,\comp{\gamma}_s,p-2+u_s+\delta)_p,
\end{align*}
with $\wt(e \star d) = (r+s+2)(p-1)+ (r+1)(p-1) + 2 > m(p - 1)+ \wt(e)$.

Finally, let us consider the case $(u_s,\delta) = (1,1)$. By \Cref{lem:technical}, there exists an index $j$ in $\{1,\dots,s\}$ such that $\gamma_j \notin\{0,1,p-2,p-1\}$. By choosing $z = j$, we obtain
\begin{align*}
e \star d = [&\overbrace{p-1,\dots,p-1}^r,\hphantom{p - }\,\;t'+1,\gamma_1,\dots,\gamma_{z-1},\gamma_z + 2, \gamma_{z+1},\dots,\gamma_s, p-u_s-3-\delta,\\ &\,p-1,\dots,p-1,p-t'+1, \comp{\gamma}_1,\dots,\comp{\gamma}_{z-1},\comp{\gamma}_{z}+2,\comp{\gamma}_{z+1},\dots,\comp{\gamma}_s,p-4+u_s+\delta]_p \\
= (&p-1,\dots,p-1,\phantom{p - }\,\;t'+1,\gamma_1,\dots,\gamma_{z-1},\gamma_z + 2, \gamma_{z+1},\dots,\gamma_s, p-5,\\ &p-1,\dots,p-1,p-t'+1, \comp{\gamma}_1,\dots,\comp{\gamma}_{z-1},\comp{\gamma}_{z}+2, \comp{\gamma}_{z+1},\dots,\comp{\gamma}_s,p-2)_p,
\end{align*}
with $\wt(e \star d) = (r+s+2)(p-1)+ (r+1)(p-1) + 2 > m(p - 1) + \wt(e)$.
\end{proof}

\end{document}
